%-------------------------------------------------------------------------------
\section{Introduction}
%-------------------------------------------------------------------------------

% - 第三方在线支付在 Web / Mobile / App 中的广泛使用
% - 第三方支付协议复杂、跨信道、跨参与方
% - 现实世界频繁出现“支付成功但订单异常”等安全事故

% #### 1.2 挑战

% 突出 **你论文独有的挑战**：

% - 支付协议不是单一协议，而是：
%   - 多支付商（微信/支付宝）
%   - 多场景（PC扫码 / H5 / App）
%   - 多信道（前端 / 后端 / 跳转信道）
% - 用户“很强”，但**收款账户与商家品牌名分离**
% - 商家前后端和跳转信道可能不安全

% 👉 这一步是为你后面的**安全假设**做铺垫

% #### 1.3 方法与贡献（Contributions）

% **明确列出贡献（USENIX 强烈偏好）：**

% > Our contributions are as follows:

% 1. **Systematic protocol analysis** of six widely deployed third-party payment protocols (2 providers × 3 scenarios).
% 2. **A layered and reusable security framework** that decomposes payment security into:
%    - Order Consistency
%    - Payment Consistency
%    - User Authorization
% 3. **Formal modeling and verification** of all protocols under realistic adversarial assumptions using TAMARIN.
% 4. **Three attack patterns** that violate the proposed security properties, affecting both payment models.
% 5. **Real-world validation and large-scale measurement** over xxxx apps/games, revealing ~20% vulnerable implementations.
% 6. **Responsible disclosure and confirmation** by WeChat Pay and Alipay.

% > 这一段如果写好，是“中稿的第一道门槛”
\clearpage